\subsection{Four-generator radial models}
\label{sec:exceptional-four-generator}

\subsubsection{Complexified Jordan norms}
\label{exc-jordan-models}

The sextic \(D\) is now independent: \(D=(\Gamma(S,S)-4RS)/12\).
In the preceding trace family, \(D=F^2\).


Let \(\mathbb K_d=\R,\mathbb C,\mathbb H,\mathbb O\) for
\(d=1,2,4,8\), respectively, and set
\(J_d=\operatorname{Herm}_3(\mathbb K_d)\), writing \(J=J_d\) locally.
Its Jordan product is \(x\bullet y=(xy+yx)/2\), its unit is \(e\),
and its Euclidean inner product is \(\tr(x\bullet y)\).
Thus \((J_d)_0=e^\perp=\{x\in J_d:\tr x=0\}\) denotes its trace-free part.
On the real vector space
underlying \(V=J\otimes_{\R}\mathbb C\), use the real part of
the Hermitian extension of this inner product. Thus
\(|x+iy|^2=|x|^2+|y|^2\) and \(\dim_{\R}V=6(d+1)\).
Let \(N\) be the complexification
of the Jordan determinant. A Jordan frame consists of three
idempotents \(c_i\) with \(c_i\bullet c_j=\delta_{ij}c_i\),
\(\tr c_i=1\), and \(\sum_i c_i=e\). The normalization is
\[
 N(z_1c_1+z_2c_2+z_3c_3)=z_1z_2z_3,
 \qquad |z_1c_1+z_2c_2+z_3c_3|^2=\sum_i|z_i|^2.
\]
These are the complexified rank-three models in
\cite[Section~6.11.6(iv)]{NadirashviliTkachevVladut14}; the polynomial
there is \(2\operatorname{Re}N\), with the same zero set.
Their REC characteristic dimensions are \((n_1,n_2)=(1,3d+2)\).

The determinant has the coordinate formula
\begin{equation}\label{exc-norm-coordinates}
 z=\begin{pmatrix}
 a&\bar u&\bar v\\ u&b&w\\ v&\bar w&h
 \end{pmatrix},\qquad
 N(z)=abh-a\,n(w)-b\,n(v)-h\,n(u)
       +2\operatorname{Sc}((\bar u w)v).
\end{equation}
Here \(a,b,h\in\mathbb C\) and
\(u,v,w\in\mathbb K_d\otimes_{\R}\mathbb C\);
the bar is algebra conjugation fixing complex scalars,
\(n(u)=u\bar u\), and \(\operatorname{Sc}\) is the complex-linear
extension of the scalar part. Thus \(n\) is a complex quadratic
form, distinct from the Hermitian norm \(|u|^2\).

The quadratic adjoint \(z^\#\) satisfies
\((dN)_z(v)=\tr(z^\#\bullet v)\), with the complex bilinear
trace form. The rank-three identities are
\begin{equation}\label{exc-adjoint}
 z\bullet z^\#=N(z)e,\qquad
 (z^\#)^\#=N(z)z,\qquad
 \tr(z\bullet z^\#)=3N(z).
\end{equation}
They extend from \(J\) to \(V\) as polynomial identities; see
\cite[Section~3.2.7]{NadirashviliTkachevVladut14} for the Jordan algebra
background. Put
\begin{equation}\label{exc-invariants}
 F=\operatorname{Re}N,\qquad R=|z|^2,\qquad
 S=|z^\#|^2,\qquad D=|N(z)|^2.
\end{equation}
Holomorphicity gives \(\Delta F=0\) and \(|\nabla F|^2=S\).
The Jordan product on \(J\) is distinct from the metrized product
obtained by polarizing \(F\) on \(V\).

\begin{theorem}[The smallest closure of the complexified models]
\label{thm:exc-jordan-closure}
For each \(d\in\{1,2,4,8\}\), the four polynomials in
\eqref{exc-invariants} are algebraically independent and
\[
 \cA_R(F)=\R[R,F,S,D].
\]
Their degrees are \(2,3,4,6\); in particular,
\((\cA_R(F))_2=\R R\) and \((\cA_R(F))_3=\R F\).
The Laplacian identities are
\begin{equation}\label{exc-jordan-laplacians}
 \Delta R=12(d+1),\qquad \Delta F=0,\qquad
 \Delta S=4(d+2)R,\qquad \Delta D=4S.
\end{equation}
The gradient pairings, apart from those determined by
\(\Gamma(R,P)=2\deg(P)P\), are
\begin{equation}\label{exc-gradient-table}
\begin{aligned}
 \Gamma(F,F)&=S,& \Gamma(F,S)&=4RF,& \Gamma(F,D)&=2SF,\\
 \Gamma(S,S)&=4RS+12D,& \Gamma(S,D)&=8RD,&
 \Gamma(D,D)&=4SD.
\end{aligned}
\end{equation}
In particular, \(\cL(F)=-2RF\), so the regular zero set is minimal.
Throughout \(V\),
\begin{equation}\label{exc-global-inequalities}
 R^2\geq3S,\qquad RS\geq9D,\qquad S^2\geq3RD,
 \qquad 0\leq F^2\leq D.
\end{equation}
\end{theorem}

\begin{proof}
\emph{Differential identities.}
For real \(a\in J\), \(T_a:z\mapsto a\bullet z\)
is self-adjoint for the Hermitian trace metric, and
\eqref{exc-adjoint} gives \((dN)_z(T_az)=(\tr a)N(z)\).
Consequently \(\exp(itT_a)\) is unitary and multiplies \(N\)
by \(\exp(it\tr a)\). Real Jordan automorphisms are also
unitary and preserve \(N\). These transformations preserve
\(R,S,D\) and transport identities involving \(N\) with its
phase factor.

The linearized Jordan identity implies that the commutators
\([T_a,T_b]\) are derivations, whose exponentials are automorphisms.
Take the standard diagonal frame \((c_i)\). At
\(x=\sum_i r_ic_i\), with distinct positive \(r_i\), the
Peirce decomposition \cite[Section~3.2.6]{NadirashviliTkachevVladut14} is
\[
 J=\bigoplus_i\R c_i\ \oplus\!\bigoplus_{i<j}J_{ij},
 \qquad \dim_{\R}J_{ij}=d.
\]
Here \(J_{ij}\) consists of matrices supported in the \((i,j)\)
and \((j,i)\) entries, and
\(c_i\bullet a=c_j\bullet a=a/2\) for \(a\in J_{ij}\).
Varying the \(r_i\) supplies the real diagonal directions, while
the derivations and unitary infinitesimal transformations give
\[
 [T_{c_i},T_a]x=\tfrac14(r_j-r_i)a,\qquad
 iT_ax=\tfrac12(r_i+r_j)ia,\qquad iT_{c_i}x=ir_ic_i.
\]
These directions span the underlying real space of \(V\).
The map obtained by varying the diagonal entries and the transformation
parameters is therefore a submersion at \(x\), so its image contains
a real open subset of \(V\).

On the diagonal, with \(t_i=r_i^2\),
\(R=e_1(t), S=e_2(t), D=e_3(t)\), where
\(e_1(t)=\sum_i t_i\), \(e_2(t)=\sum_{i<j}t_it_j\), and
\(e_3(t)=t_1t_2t_3\). The first derivatives of
these three invariants in the off-diagonal directions vanish.
Differentiation gives the last row of \eqref{exc-gradient-table}.
Extending \(\Gamma\) complex bilinearly, diagonal differentiation
gives \(\Gamma(N,S)=4RN\), whose real part is the desired identity
for \(F\). Holomorphicity gives the remaining pairings with
\(F\). The complex Hessian of \(N\) has diagonal-block
squared norm \(2R\); on \(J_{ij}\) its block is
\(-r_k\Id\), where \(\{i,j,k\}=\{1,2,3\}\), contributing
\(dR\) in total. Therefore
\[
\begin{gathered}
 |\Hess_{\mathbb C}N|^2=(d+2)R,\qquad
 |\Hess_{\R}F|^2=2(d+2)R,\\
 \Delta S=2|\Hess_{\R}F|^2=4(d+2)R.
\end{gathered}
\]
Here \(|\Hess_{\R}\operatorname{Re}(e^{i\theta}N)|^2
=2|\Hess_{\mathbb C}N|^2\), so the real Hessian norm is
unchanged by the determinant phase.
Also \(\Delta|N|^2=4|z^\#|^2\). Phase covariance transports
these identities to the open subset just obtained, and polynomiality
then proves them on all of \(V\).

\emph{The smallest closure.}
The differential identities satisfy
\cref{prop:cubic-scalar-killing-closure}, which gives the claimed
smallest algebra. At distinct positive \(t_i\), the Jacobian
of \((e_1,e_2,e_3)\) is nonzero. Changing the overall complex
phase preserves \(R,S,D\) and changes \(F\) whenever \(N\)
is nonreal. Hence the four differentials are independent at such
a point.

\emph{Global inequalities.}
The gradient table and \cref{lem:cubic-unified-bounds} give
\eqref{exc-global-inequalities}.
\end{proof}

\subsubsection{Polynomial potentials and the Jordan minimizing range}
\label{exc-multipliers}

The following polynomial alternative also applies to the model in
\cref{exc-g2-model}.

\begin{proposition}[An algebraic sign criterion for the multiplier \(S\)]
\label{prop:exc-SF-certificate}
Let \(F\) be a harmonic cubic on a Euclidean space \(V\), set
\(R=|x|^2\) and \(S=|\nabla F|^2\), and let \(D\) be a
sextic satisfying \eqref{exc-gradient-table} and
\eqref{exc-global-inequalities}. Suppose \(\Delta S=cR\).
For \(G=SF\), one has \(\cL(G)=F\mathcal P_c\), where
\begin{equation}\label{exc-SF-certificate}
\begin{aligned}
 \mathcal P_c={}&(c-\tfrac{58}{3})RS^3
       +\tfrac{10}{3}S^2(RS-9D)\\
 &+F^2\bigl[(12c-72)RD+12cS(R^2-3S)+(36c-96)S^2\bigr].
\end{aligned}
\end{equation}
If \(c\geq58/3\), then \(G\cL(G)\geq0\), and
\(\nabla G\neq0\) wherever \(S>0\).
\end{proposition}

\begin{proof}
Set \(K=cR\) in \eqref{eq:cubic-SF-general-defect} and regroup
to obtain \eqref{exc-SF-certificate}. Every summand is nonnegative
by \eqref{exc-global-inequalities} when \(c\geq58/3\).
The gradient formula \eqref{eq:cubic-SF-general-gradient} gives
\(|\nabla G|^2\geq S^3\).
\end{proof}

\begin{theorem}[Area minimization for the complexified Jordan models]
\label{thm:exc-jordan-minimizing}
For the four models in \cref{exc-jordan-models},
\(\{F<0\}\) is locally perimeter minimizing if and only if
\(d\in\{2,4,8\}\). In all three minimizing dimensions,
\(G=F\sqrt{RS}\) generates a subcalibration. For \(d=4,8\),
the polynomial potential \(SF\) does so as well.
\end{theorem}

\begin{proof}
For \(d=2,4,8\), the dimension and trace coefficient are
\[
 (N,K/R)=(18,16),\ (30,24),\ (54,40),\qquad K=\Delta S.
\]
Since \(\tr L_x^2=K/2\), each pair satisfies
\cref{cor:cubic-unified-trace}; this proves the minimizing assertions
for the common potential. For \(d=4,8\),
\cref{prop:exc-SF-certificate} additionally gives the polynomial
potential \(SF\).

The singular set is the Jordan rank-at-most-one cone
\[
 \Sing\{F=0\}=\Sigma=\{z\in V:z^\#=0\}.
\]
To describe its nonzero stratum, take
\(a\neq0\) in \eqref{exc-norm-coordinates} and set
\[
 B=b-\frac{n(u)}a,\quad C=h-\frac{n(v)}a,\quad
 W=w-\frac{u\bar v}a.
\]
The determinant formula gives
\begin{equation}\label{exc-schur}
 N(z)=a\bigl(BC-n(W)\bigr).
\end{equation}
This identity also holds for \(d=8\). Indeed, multiplicativity
of \(n\) and associativity of the scalar part give
\[
 n\!\left(w-\frac{u\bar v}{a}\right)
 =n(w)-\frac{2}{a}\operatorname{Sc}((\bar u w)v)
       +\frac{n(u)n(v)}{a^2},
\]
which proves \eqref{exc-schur} after expanding \(aBC\).
The coordinate change is holomorphically invertible, so
\(dN=0\) is exactly \(B=C=W=0\) in this chart. Hence
these points are geometrically singular: the transverse real
quadratic form \(\operatorname{Re}(a(BC-n(W)))\) is nondegenerate
and indefinite. Moreover,
\(\Sigma\setminus\{0\}\) is smooth of complex dimension
\(1+2d\), and real codimension \(2(d+2)\). To see that these
charts cover the nonzero set, recall that the spectral theorem makes
primitive idempotents span \(J\), and that real Jordan automorphisms
act transitively on them
\cite[Theorem~III.1.2 and Corollary~IV.2.7]{FarautKoranyi94}.
Some primitive idempotent therefore has nonzero trace pairing with
any nonzero \(z\), and an automorphism takes it to \(c_1\).

For \(d=4,8\), the codimension bound and
\eqref{eq:cubic-SF-general-gradient} allow
\cref{cor:gradient-subcalibration} to be applied to \(SF\).
Harmonicity identifies the boundary-current support with \(\{F=0\}\)
by \cref{rem:polynomial-finite-perimeter}.

For \(d=1\), take \(p\in\Sigma\setminus\{0\}\) and put
\(H=\Hess F(p)\). The chart \eqref{exc-schur} shows that
\(H\) has real rank \(2(d+2)\).
\Cref{prop:cubic-general-singular-spectrum}, with \(q=2R\), gives
\(H^3=R(p)H\). Since \(\tr H=0\), its nonzero eigenvalues
are \(\pm\sqrt{R(p)}\), each of multiplicity \(d+2\).
Thus the tangent cone, with the induced Euclidean metric, is
\begin{equation}\label{exc-jordan-tangent}
 C_{d+1,d+1}\times\R^{4d+2},\qquad
 C_{q,q}=\{(x,y)\in\R^{q+1}\oplus\R^{q+1}:|x|=|y|\}.
\end{equation}
For \(d=1\), the normal factor \(C_{2,2}\) is unstable by the
Hardy criterion in \cref{sec:tangent-instability-transfer}.
\Cref{lem:tangent-instability-transfer} therefore rules out local
perimeter minimization.
\end{proof}

For \(d=2\), complexification identifies \(V\) isometrically
with \(M_3(\mathbb C)\) with its Frobenius metric, and \(N\)
with the determinant. In the notation of \cref{sec:complex-determinants},
\(R=E_1\), \(S=E_2\), and \(D=E_3\). Thus both the closure
table and the potential above are the three-by-three cases of the
construction there.

\subsubsection{The \texorpdfstring{$G_2$}{G2} trilinear model}
\label{exc-g2-model}

Let \(\times\) be the standard cross product on
\(\operatorname{Im}\mathbb O=\R^7\), and set
\[
 F(a,b,c)=\langle a\times b,c\rangle,
 \qquad (a,b,c)\in(\R^7)^3.
\]
Equivalently, \(F(a,b,c)=\varphi(a,b,c)\), where \(\varphi\)
is the associative calibration on \(\R^7\) of Harvey and Lawson
\cite[Section~IV.1.A]{HarveyLawson82a}.
This is the classical model in
\cite[Section~6.4, equation~(6.26), Proposition~6.3]{TkachevClassification10},
whose defining polynomial \(\operatorname{Re}(a(bc))\) is
\(-F\). Up to similarity, its subtracted Jordan model is the cubic
norm restricted to the trace-orthogonal complement of
\(\operatorname{Herm}_3(\R)\) in \(\operatorname{Herm}_3(\mathbb O)\);
see \cite[Proposition~6.4.4 and Table~6.4.1]{NadirashviliTkachevVladut14}.
Let \(A\) be the real \(7\times3\) matrix with
columns \(a,b,c\), and put
\[
 B=A^tA,\qquad R=\tr B,\qquad S=e_2(B),\qquad D=\det B,
\]
where \(e_2(B)\) is the second elementary symmetric polynomial
of the eigenvalues of \(B\).

The compact real group \(G_2=\operatorname{Aut}(\mathbb O)\)
acts diagonally on \(V=(\R^7)^3\). Schwarz's invariant-generation
theorem \cite[Theorem~4 and equations~(5.1)--(5.3)]{Schwarz83} gives
\[
 \R[V]^{G_2}=\R[B_{ij},F].
\]
The commuting \(SO(3)\) action on the columns fixes \(F\)
and conjugates \(B\). Averaging polynomial expressions over
\(SO(3)\), whose conjugation invariants on symmetric matrices
are generated by the elementary symmetric polynomials, therefore gives
\[
 \R[V]^{G_2\times SO(3)}=\R[R,F,S,D].
\]
This invariant ring equals \(\cA_R(F)\):

\begin{theorem}[The trilinear model]
\label{thm:exc-g2}
The four polynomials \(R,F,S,D\) are algebraically independent,
and \(\cA_R(F)=\R[R,F,S,D]\). In particular,
\((\cA_R(F))_2=\R R\) and \((\cA_R(F))_3=\R F\).
They satisfy
\eqref{exc-gradient-table} and \eqref{exc-global-inequalities},
with
\begin{equation}\label{exc-g2-laplacians}
 \Delta R=42,\qquad \Delta F=0,\qquad
 \Delta S=24R,\qquad \Delta D=10S.
\end{equation}
Both \(G=F\sqrt{RS}\) and the polynomial potential \(SF\)
give subcalibrations of \(\{F<0\}\subset\R^{21}\), which is
locally perimeter minimizing.
\end{theorem}

\begin{proof}
The cross-product identity
\(|a\times b|^2=|a|^2|b|^2-\langle a,b\rangle^2\)
gives \(|\nabla F|^2=S\), while trilinearity gives
\(\Delta F=0\). Alternation gives
\(\langle\nabla_{a_i}F,a_j\rangle=\delta_{ij}F\), with
\(a_i\) denoting the columns. In matrix notation,
\[
 \nabla_A S=2A(RI-B),\qquad
 \nabla_A D=2A\operatorname{adj}B.
\]
These formulas and the three-by-three Newton identities give
\eqref{exc-gradient-table}. Direct differentiation of the Gram
invariants gives \(\Delta S=24R\) and \(\Delta D=10S\).
The global inequalities \eqref{exc-global-inequalities} follow
from \cref{lem:cubic-unified-bounds}.

The smallest-algebra assertion now follows from
\cref{prop:cubic-scalar-killing-closure}.
For independence,
choose orthonormal \(e_1,e_2\), put \(e_3=e_1\times e_2\),
and choose a unit \(e_4\) orthogonal to their span. At
\[
 a=r_1e_1,\qquad b=r_2e_2,\qquad
 c=r_3(\cos\theta\,e_3+\sin\theta\,e_4),
\]
the three invariants \(R,S,D\) are the elementary symmetric
functions of \(r_1^2,r_2^2,r_3^2\), while
\(F=r_1r_2r_3\cos\theta\). Distinct positive \(r_i^2\)
and \(0<\theta<\pi\) give four independent differentials.

Here \(N=21\) and \(K=24R\), so
\cref{cor:cubic-unified-trace} gives the common potential.
For the polynomial potential, \cref{prop:exc-SF-certificate} applies
with \(c=24\). The singular set of the zero cone is
\[
 \Sing\{F=0\}=\Sigma=\{S=0\}
 =\{A\in M_{7,3}(\R):\operatorname{rank}A\leq1\}.
\]
At a nonzero rank-one matrix, a column rotation reduces to
\((a,0,0)\), whose transverse leading form
\(\langle a\times\delta b,\delta c\rangle\) is indefinite of rank twelve.
The nonzero stratum has dimension
\(7+3-1=9\), hence codimension \(12\). The gradient bound
\eqref{eq:cubic-SF-general-gradient} and \cref{cor:gradient-subcalibration}
therefore give the same minimizing phase.
\end{proof}

\begin{remark}
For both families, Bochner's identity gives
\(2|L_x|_{\mathrm{HS}}^2=c|x|^2\), where \(c>0\). Hence
\(L_v=0\) implies \(v=0\), and their smallest closures have the intrinsic
degree pattern \((3)\) of \cref{thm:cubic-one-dimensional-degree-three}.
For the \(G_2\) model this differs from its trilinear coordinate
presentation: the larger invariant algebra
\(\R[V]^{G_2}\) contains their separate radii, whereas the
smallest algebra is also invariant under column rotations and
does not contain \(|a|^2,|b|^2,|c|^2\).
\end{remark}
